% =============================================================================
% 1.3 主要研究内容
% 草稿版本：v0.1 (2026-05-05)
% 字数：约 1500 字
% 风格规则：v0.2 风格（无 \textbf / \textit / 引例括号 / 双引号 / 破折号）
% 引用：暂无新增引用
% 图表：
%   - 表 \reftab{tab:contributions}    五项创新贡献与对应实证概览
% 依赖宏包：booktabs / tabularx（主稿已加载）
% =============================================================================

\subsection{主要研究内容}
\label{sec:research-content}

针对 \ref{sec:related-work} 节归纳出的三方面研究不足，本文把研究
内容收敛为三个核心研究问题，并据此提出五项创新贡献。

\subsubsection*{1.3.1 \quad 三个研究问题}

研究问题 Q1 是气象任务指令解析问题。给定用户的自然语言查询,
如何把它精确映射为结构化的检索意图与可执行的 API 调用参数。
该问题涉及命名实体识别、相对时间消歧、参数枚举映射、缺失参数
推理与对抗鲁棒性五类子任务。

研究问题 Q2 是气象数据的异构特征分析问题。气象数据在数据源、
数据粒度、数据要素三个层面均存在显著的异构性，如何在不损失原始
数据精度的前提下把多源数据统一为大语言模型可消费的标准接口与
语义化语料，是气象 Agent 中段流水线必须解决的关键问题。

研究问题 Q3 是大语言模型在数值计算与逻辑推理中的幻觉问题。气象
数据的统计分析与决策报告生成对数值精度与权威引用要求严格，如何
通过架构层面的设计避免大语言模型在数值边界判定、统计计算、标准
引用等高风险任务上的失败，是本文需要解决的最具挑战性的问题。

\subsubsection*{1.3.2 \quad 五项创新贡献}

围绕上述三个研究问题，本文给出以下五项创新贡献，与对应的章节
位置以及评测出口共同列于 \reftab{tab:contributions}。

\begin{table}[!htbp]
    \centering
    \caption{本文五项创新贡献与对应章节及评测概览}
    \label{tab:contributions}
    \song\fontsize{9pt}{12pt}\selectfont
    \renewcommand{\arraystretch}{0.95}
    \begin{tabular}{lll}
        \toprule
        贡献 & 章节 & 关键实证 \\
        \midrule
        ReAct 单智能体气象 Agent 架构 & \ref{sec:react-paradigm} & 11 个工具集成、单进程多轮对话 \\
        工具学习驱动的异构数据检索 & \ref{chap:retrieval} & 端到端 91.4\% / 角色 100\% \\
        双通路 RAG 与语义桥接 & \ref{sec:semantic-bridge} & Top-1 85.0\% / \texttt{grade\_\bk{}id} 100\% \\
        代码生成与受限沙箱执行 & \ref{sec:code-execution} & 数值准确率 +22.2 pp \\
        多层评测体系（5 套 214 用例） & \ref{chap:evaluation} & 端到端通过率 +13.4 pp \\
        \bottomrule
    \end{tabular}
\end{table}

第一项贡献是基于 ReAct 范式的气象 Agent 单智能体架构。本文采用
单智能体加显式工具学习而非多智能体协作的工程取舍，把场景编排、
出处强制、出行场景多地点对齐等约束下沉到 System Prompt，避免
了多智能体之间的状态同步与额外延迟。该架构集成 9 个气象 API
工具加 2 个 RAG 工具共 11 个工具单元，并通过 LangGraph 的
\texttt{InMemorySaver} 与 \texttt{thread\_\bk{}id} 机制实现多轮对话
状态保持。

第二项贡献是工具学习驱动的异构数据检索机制。本文为和风天气与
Open-Meteo 两类异构 API 设计统一的 Function Calling Schema,
并通过意图识别与参数补全双阶段管线把自然语言查询映射为结构化
\texttt{WeatherIntent} 对象。在 35 条评测集上端到端通过率达到
91.4\%，多地点角色识别准确率达到 100\%，证明了分层架构在容错
能力上的天然优势。

第三项贡献是双通路 RAG 架构。通路 A 把向量检索与 BM25 词法检索
以可调权重融合，承担概念性查询的高召回任务；通路 B 由 5 类确定性
分级器加 \texttt{grade\_\bk{}id} 硬链接构成，承担数值到等级再到标准
引用的精确归属任务。在 60 条检索评测集与 71 条桥接评测集上,
通路 A 的 Top-1 命中率达到 85.0\%、MRR 达到 0.908，通路 B 的
\texttt{grade\_\bk{}id} 准确率与权威引用真实性同时达到 100\%，与
大语言模型自由发挥的基线在严格命中率上形成 74.6 个百分点差距。

第四项贡献是基于代码生成与受限沙箱的确定性统计分析机制。本文
设计的轻量沙箱方案由大语言模型生成 \texttt{compute(data)} 顶层
函数，在白名单 builtins 加白名单模块加墙钟超时加结构化错误捕获
四重约束下执行。在 18 条覆盖 6 类统计模式的评测集上，代码执行档
数值准确率达到 88.9\%，比大语言模型自由直算高出 22.2 个百分点,
失败模式分析进一步揭示了大语言模型在伴随状态维护与离散计数任务
上的内在边界。

第五项贡献是覆盖模块层与集成层的多层评测体系。本文构建了 5 份
独立评测集共 214 条用例，包含 4 个模块层评测加 1 个 30 条端到端
集成层评测。集成层评测以最终用户体验视角验证整套系统的端到端
价值，full 档端到端通过率为 86.7\%，相对 ablated 去 RAG 对照档
提升 13.4 个百分点；权威标准引用率从 0\% 提升到 26.7\%；
decision\_support 类用例通过率从 50\% 提升到 100\%。同时该评测
也揭示了从模块层到集成层的衰减效应，构成 \ref{sec:future-work}
节未来工作的天然指向。

\subsubsection*{1.3.3 \quad 研究方法}

本文的研究方法以工程实现为主线，沿用 ReAct 范式的设计原则，并
辅以系统化实证。具体研究路径分为四个阶段。第一阶段是文献调研与
方案设计，通过研读 ReAct、CoT、RAG、工具学习等核心工作 \cite{ref13,
ref18, ref12, ref21, ref28} 形成总体方案。第二阶段是模块实现,
按 Q1 接入侧、Q2 中段流水线、Q3 解读侧的顺序依次落地。第三阶段
是模块层评测，分别针对意图识别、检索、桥接、代码执行 4 个模块
独立设计评测集与消融对照。第四阶段是集成层评测与论文撰写，
通过 30 条真实场景用例的 full 档与 ablated 档对照验证整套系统的
端到端价值，并据此完成论文核心论证闭环。

